\section{Motivation for Migrating from CSV to PostgreSQL}
\begin{tcolorbox}[title=Note]
CSV is unsuitable for server environments because: (1) it does not support concurrent access, (2) data is lost when a container restarts, and (3) conditional queries are inefficient. PostgreSQL resolves all of these issues while providing full ACID compliance.
\end{tcolorbox}

\section{Database Schema}

\begin{lstlisting}[language=sql,caption={PostgreSQL schema for the Fashion Agent}]
-- Main product table
CREATE TABLE IF NOT EXISTS fashion_items (
    id            SERIAL PRIMARY KEY,
    image_id      TEXT UNIQUE NOT NULL,
    label         TEXT,
    color         TEXT,
    caption       TEXT,
    source        TEXT DEFAULT 'kaggle',
    image_path    TEXT,
    created_at    TIMESTAMPTZ DEFAULT NOW()
);

-- User session table
CREATE TABLE IF NOT EXISTS sessions (
    session_id          TEXT PRIMARY KEY,
    created_at          TIMESTAMPTZ DEFAULT NOW(),
    query_history       JSONB DEFAULT '[]',
    liked_items         JSONB DEFAULT '[]',
    gender              TEXT DEFAULT 'unisex',
    gender_hint_enabled BOOLEAN DEFAULT TRUE,
    year_of_birth       INTEGER
);

-- Conversation message table
CREATE TABLE IF NOT EXISTS session_messages (
    id            SERIAL PRIMARY KEY,
    session_id    TEXT REFERENCES sessions(session_id),
    role          TEXT NOT NULL,  -- 'user' or 'assistant'
    content       TEXT NOT NULL,
    created_at    TIMESTAMPTZ DEFAULT NOW()
);
\end{lstlisting}

\section{Data Access Layer}

\begin{lstlisting}[caption={PostgreSQL Data Access Layer --- agent/memory.py}]
import psycopg2

def _get_conn():
    """Connect to PostgreSQL from environment variables."""
    return psycopg2.connect(
        host=os.getenv("PGHOST", "localhost"),
        port=int(os.getenv("PGPORT", "5432")),
        dbname=os.getenv("PGDATABASE", "fashion_rag"),
        user=os.getenv("PGUSER", "fashion_user"),
        password=os.getenv("PGPASSWORD", ""),
    )

def create_session() -> str:
    """Create a new session and return its UUID."""
    sid = str(uuid.uuid4())
    with _get_conn() as conn:
        with conn.cursor() as cur:
            cur.execute(
                "INSERT INTO sessions (session_id) VALUES (%s)",
                (sid,))
        conn.commit()
    return sid

def get_preferences(session_id: str) -> dict:
    """Extract accumulated preferences from query history."""
    with _get_conn() as conn:
        with conn.cursor() as cur:
            cur.execute(
                "SELECT query_history FROM sessions "
                "WHERE session_id = %s", (session_id,))
            row = cur.fetchone()
    if not row or not row[0]:
        return {}
    # Analyze filters from history to derive preferred colors/categories
    history = row[0] if isinstance(row[0], list) else []
    colors, categories = Counter(), Counter()
    for entry in history:
        filters = entry.get("filters", {})
        if filters.get("color"):
            colors[filters["color"]] += 1
        if filters.get("category"):
            categories[filters["category"]] += 1
    return {
        "preferred_colors": [c for c, _ in colors.most_common(3)],
        "preferred_categories":
            [c for c, _ in categories.most_common(3)],
    }
\end{lstlisting}

\textbf{Task breakdown:}
\begin{itemize}
    \item \textbf{Environment-based connection}: Connects to PostgreSQL via environment variables, suitable for Docker deployment.
    \item \textbf{JSONB}: Stores liked items and query history flexibly without requiring additional tables.
    \item \textbf{Preference extraction}: Automatically aggregates the top-3 colors and categories from query history.
\end{itemize}

% ===========================================================================
% CHAPTER 5: FASHION AGENT
% ===========================================================================
